% Data note 2026-07-30: this file contains EXPECTED values for the per-image prompt-tightening/oracle-description data setting.
% It should not be mixed with the current Ours (unnamed) prototype-reference data setting unless explicitly revived.
零样本异常检测（ZSAD）借助 CLIP 的通用异常 prompt，在无需目标类别训练数据的情况下定位缺陷。这类文本表述具备良好的跨类别泛化能力，却只能在类别无关的宽泛语义上刻画“异常”，无法描述当前图像中异常的实际形态；由此异常侧语义过宽，使得显著偏离正常模板的正常结构（如规则纹理、花纹）也被划入异常一侧，产生误响应。我们观察到，若在通用 prompt 之外引入面向当前图像的具体异常描述，异常侧的语义得以收紧，模型响应随之向真实异常区域集中。然而，收紧这一侧所需要的位置与形态信息无法由通用 prompt 预先写定，须在测试阶段逐图像地从图像自身获取。为此，本文提出一种高频驱动的逐图 prompt 自适应方法：以 CLIP patch 特征的高频分量所刻画的局部结构偏离作为无监督条件下的可计算依据，在测试时逐图收紧异常侧文本表示，使其对齐到当前图像的实际异常之上。整个过程冻结 CLIP、无需训练与反向传播。在 MVTec AD、VisA、MPDD、BTAD 与 DTD-Synthetic 五个基准上，本方法在像素级定位指标（pixel AUPRO）上较通用 prompt 基线取得一致提升，同时不在正常图像上引入额外误报。（本文实验数值为预期目标值，用于展示方法结构，待真实实验回填。）
